\chapter{Autocallables}\label{ch:dv:autocallables}

In January 2024 the CSI 500 and CSI 1000 indices of mainland Chinese shares fell sharply. Structured notes known as
\omterm{def:dv:autocallables:snowball}{snowballs}, with knock-in levels between 70\% and 80\% of their starting points, were knocked in as the indices fell. The
notes paid coupons of 10\% to 20\% a year for as long as the indices stayed within their range. Once knocked in, they turned into positions that lost with the index. The dealers who had issued them held
hedges that changed abruptly at those levels. The same kind of product, the \omterm{def:dv:autocallables:autocallable}{autocallable}, is sold to retail investors in
Korea and Japan. Korean investors held 19.3 trillion won of notes linked to the Hang Seng China Enterprises Index in
November 2023, 10.2 trillion of them due in the first half of 2024; by 22 January 2024 those sold by five commercial banks had
already lost 229.6 billion won, and the losses have been studied in their own right. This chapter takes the product apart. It covers the payoff and the model it needs, its risks for the
issuer, the hedging flows it creates, and what those flows have done to the markets they hedge in.

\section{Payoff anatomy}

\begin{definition}[Autocallable]\label{def:dv:autocallables:autocallable}
An \emph{autocallable}\index{autocallable} is a note that redeems early, at par plus a coupon, on the first observation
date on which the underlying is at or above a trigger level. If it is never called, it repays par at maturity, unless a
\omterm{def:dv:autocallables:protection}{protection barrier} has been breached, in which case it repays par times the underlying's performance.
\end{definition}

\begin{definition}[Autocall trigger]\label{def:dv:autocallables:trigger}
The \emph{autocall trigger}\index{autocall trigger} is the level, usually 100\% of the initial fixing and sometimes
stepping down over time, at or above which the note is redeemed early on an observation date.
\end{definition}

\begin{definition}[Coupon barrier]\label{def:dv:autocallables:couponbarrier}
The \emph{coupon barrier}\index{coupon barrier} of a Phoenix \omterm{def:dv:autocallables:autocallable}{autocallable} is the level at or above which a periodic coupon is
paid on an observation date, whether or not the note is called.
\end{definition}

\begin{definition}[Memory coupon]\label{def:dv:autocallables:memory}
A \emph{memory coupon}\index{memory coupon} is a Phoenix coupon that, if missed because the underlying was below the \omterm{def:dv:autocallables:couponbarrier}{coupon
barrier}, is paid later on the first observation date on which the barrier is met again.
\end{definition}

\begin{definition}[Protection barrier]\label{def:dv:autocallables:protection}
The \emph{protection barrier}\index{protection barrier} (knock-in level) is the level below which the investor's capital is
no longer protected. It is observed at maturity only (European) or continuously or daily (American), and it is set well below
the initial fixing: 60\% in the chapter's example, 70\% to 80\% in the Chinese \omterm{def:dv:autocallables:snowball}{snowballs} of 2024.
\end{definition}

Taken apart, an \omterm{def:dv:autocallables:autocallable}{autocallable} is a bond plus a set of options, and the investor is short most of them. The investor has sold
a down-and-in put struck at the initial level (chapter 15's knock-in): below the \omterm{def:dv:autocallables:protection}{protection barrier} the note loses with the
underlying. The investor has bought a strip of \omterm{def:dv:barriers-and-digitals:digital}{digital options} that pay the coupons, and sold the upside above par, because
the note is called away when the underlying rises. The coupon is financed by the knock-in put premium and by the early
call, which cuts the note's life short in exactly the scenarios in which the investor would have liked to keep it.

\begin{definition}[Snowball]\label{def:dv:autocallables:snowball}
A \emph{snowball}\index{snowball} is an \omterm{def:dv:autocallables:autocallable}{autocallable}, common in China, whose coupon accrues at an annual rate and is paid in a
lump when the note knocks out (autocalls) on a monthly observation, or at maturity if it neither knocked out nor knocked in;
the knock-in is monitored daily, and after a knock-in the investor bears the underlying's loss at maturity.
\end{definition}

The chapter's working example is a three-year Phoenix on the index of chapter 9's surface. It is observed quarterly and
callable from the first quarter at 100\%. Its \omterm{def:dv:autocallables:memory}{memory coupon} is paid at or above 70\%, and its \omterm{def:dv:autocallables:protection}{protection barrier} is observed
at maturity at 60\%. Rates are zero, so that the whole coupon is paid for by options.

\begin{example}[Where the Phoenix goes]\label{ex:dv:autocallables:outcomes}
Under chapter 9's \omterm{def:dv:local-volatility:model}{local volatility}, the note is called at the first observation with probability 57.6\%, and within the first
year with probability 80.5\%. It reaches maturity unbroken with probability 5.7\%, and knocked in with probability 6.3\%
(\cref{fig:dv:autocallables:outcomes}). The fair quarterly coupon, the one that prices the note at par, is 1.55\%, or 6.18\% a
year. Without coupons the note is worth 95.99, and each point of quarterly coupon adds 2.60.
\end{example}

\begin{omfigure}
\begin{tikzpicture}
\begin{axis}[width=12cm,height=5cm,ybar,bar width=10pt,
  xlabel={outcome},ylabel={probability (\%)},
  tick label style={font=\small},label style={font=\small},
  ymin=0,ymax=68,xmin=-0.7,xmax=12.7,
  xtick=data,xticklabels from table={figdata/derivatives/18-autocallables/outcomes.csv}{label},
  x tick label style={font=\scriptsize},
  nodes near coords,nodes near coords style={font=\tiny,/pgf/number format/fixed,/pgf/number format/precision=1},
  grid=major,grid style={gray!25},table/col sep=comma]
\addplot[omThm,fill=omThm!40] table[x=i,y=prob]{figdata/derivatives/18-autocallables/outcomes.csv};
\end{axis}
\end{tikzpicture}

\omcaption{The fates of a three-year quarterly Phoenix (trigger 100\%, \omterm{def:dv:autocallables:couponbarrier}{coupon barrier} 70\%, protection 60\% at maturity) under
chapter 9's \omterm{def:dv:local-volatility:model}{local volatility}: called at each observation, redeemed at par at maturity, or knocked in. Most notes are called in
the first year; the investor's loss sits in the last bar. Data: the tutorial.}\label{fig:dv:autocallables:outcomes}
\end{omfigure}

\section{Pricing: which model}

The payoff depends on the path at a dozen dates and on the smile at every one of them. The coupons are digitals, so they are
priced by the skew (chapter 15). The knock-in put is a barrier deep out of the money, priced by the downside of the smile at
long expiries. And the early call is a string of forward digitals. A flat volatility gets all three wrong in the same
direction.

\begin{example}[Flat volatility against the smile]\label{ex:dv:autocallables:flat}
Priced at a flat 20.0\%, the three-year at-the-money volatility, the Phoenix with the local-volatility fair coupon is worth 102.78
instead of 100. The fair coupon at flat volatility is 3.14\% a year, half of the 6.18\% under \omterm{def:dv:local-volatility:model}{local volatility}. The skew makes
the knock-in put, which the investor has sold, more valuable, and the smile's steeper short end lowers the probability of an
early call. A desk that quoted with flat volatility would pay half the coupon it could afford. A competitor using the smile would
take every trade.
\end{example}

\omterm{def:dv:local-volatility:model}{Local volatility} is the natural starting point: it fits every vanilla, and the note's value depends mostly on terminal
distributions at the observation dates. It is not the end point. The coupons and calls are forward digitals, whose price
depends on the forward skew that \omterm{def:dv:local-volatility:model}{local volatility} flattens (chapter 16). The knock-in, monitored daily in many products, depends
on the smile's dynamics near the barrier. A careful desk prices with a local-stochastic volatility model (chapter 20), checks it against \omterm{def:dv:local-volatility:model}{local volatility}, and holds a
reserve for the difference. It adds a jump or gap reserve for daily-monitored barriers (chapter 15). Multi-underlying notes add correlation, and the \omterm{def:dv:multi-asset-options:skew}{correlation skew} of chapter 17.

\section{The risk profile: vega, skew, dividends, correlation}

\begin{example}[The Phoenix's sensitivities]\label{ex:dv:autocallables:greeks}
At the local-volatility fair coupon, a rise of one volatility point across the surface lowers the note's value by 0.35 per 100.
By expiry bucket, the change is $-0.16$ at one year, $-0.05$ at two years and $-0.14$ at three
(\cref{fig:dv:autocallables:vega}). A steeper skew (SSVI correlation from $-0.6$ to $-0.7$) costs 0.24. A dividend yield of 1\% a
year, with the surface unchanged in forward moneyness, costs 0.20. Between four points down and six points up the value
falls almost linearly, from 101.33 to 97.90.
\end{example}

The investor is short volatility, short skew and short dividends, and the issuer holds the opposite. The issuer is long \omterm{def:dv:greeks-and-the-hedging-pnl:greeks}{vega},
concentrated at the short end, where the first observations sit, and at maturity, where the knock-in put lives. It is long
skew and long dividends: a cut in expected dividends raises the forward and lowers the value of the knock-in put it owns. None
of these can be left open. The desk sells volatility and skew back to the market, in listed options and \omterm{def:dv:variance-swaps-and-volatility-derivatives:vs}{variance swaps}, and it
sells dividend exposure, in dividend futures and swaps (chapter 5). A market in which many issuers sell the same notes on the
same indices therefore has a structural supplier of long-dated volatility and a structural seller of dividends. That supply should
cheapen long-dated volatility and \omterm{def:dv:dividends-borrow-and-forwards:implied}{implied dividends}. The chapter does not measure by how much.

\begin{omfigure}
\begin{tikzpicture}
\begin{groupplot}[group style={group size=2 by 1,horizontal sep=1.6cm},
  width=6.6cm,height=5cm,
  tick label style={font=\small},label style={font=\small},grid=major,grid style={gray!25},table/col sep=comma]
\nextgroupplot[ybar,bar width=12pt,xlabel={volatility bump},ylabel={value change (per 100)},ymin=-0.4,ymax=0,xmin=-0.6,xmax=3.6,
  xtick=data,xticklabels from table={figdata/derivatives/18-autocallables/vega.csv}{pillar},
  title={\small vega by expiry, +1 point}]
\addplot[omThm,fill=omThm!40] table[x=i,y=vega]{figdata/derivatives/18-autocallables/vega.csv};
\nextgroupplot[xlabel={parallel volatility shift (points)},ylabel={note value (per 100)},xmin=-5,xmax=7,ymin=97.5,ymax=101.6,
  title={\small value against volatility}]
\addplot[omDef,very thick,mark=*,mark size=1.6pt] table[x=shift,y=price]{figdata/derivatives/18-autocallables/price_vol.csv};
\end{groupplot}
\end{tikzpicture}

\omcaption{The Phoenix at its fair coupon. Left: the change in value for one volatility point added around each pillar of the
surface, and to all. Right: value against a parallel shift of the surface. The investor is short \omterm{def:dv:greeks-and-the-hedging-pnl:greeks}{vega}, and the issuer long.
Data: the tutorial.}\label{fig:dv:autocallables:vega}
\end{omfigure}

\omterm{def:dv:autocallables:autocallable}{Autocallables} are written on single indices and on baskets. Worst-ofs on three or four shares pay higher coupons because the
investor sells more: a worst-of knock-in put, whose value rises as correlation falls (chapter 17).

\begin{example}[The same note on the worst of three]\label{ex:dv:autocallables:worstof}
Put the Phoenix on the worst of three underlyings, each with the index's \omterm{def:dv:local-volatility:model}{local volatility}. The fair coupon is 11.5\% a year at
a correlation of 0.5 and 10.4\% at 0.7. The knock-in probability is 18.1\% and 14.6\%. Twenty points of correlation are worth a
point of annual coupon, and the issuer who sold them is short correlation: it gains if correlation rises.
\end{example}

\begin{omfigure}
\begin{tikzpicture}
\begin{axis}[width=10.5cm,height=5cm,ybar,bar width=22pt,
  xlabel={model and underlying},ylabel={fair coupon (\% a year)},
  tick label style={font=\small},label style={font=\small},
  ymin=0,ymax=13,xmin=-0.6,xmax=3.6,
  xtick=data,xticklabels from table={figdata/derivatives/18-autocallables/coupons.csv}{label},
  nodes near coords,nodes near coords style={font=\footnotesize,/pgf/number format/fixed,/pgf/number format/precision=1},
  grid=major,grid style={gray!25},table/col sep=comma]
\addplot[omDef,fill=omDef!35] table[x=i,y=annual]{figdata/derivatives/18-autocallables/coupons.csv};
\end{axis}
\end{tikzpicture}

\omcaption{The fair annual coupon of the same three-year Phoenix: at a flat volatility, under \omterm{def:dv:local-volatility:model}{local volatility}, and on the worst
of three underlyings at correlations of 0.7 and 0.5. The smile doubles the coupon, and the worst-of nearly doubles it again. Data:
the tutorial.}\label{fig:dv:autocallables:coupons}
\end{omfigure}

\section{Issuer hedging flows}

The issuer hedges the note's delta in the underlying, and the delta of an \omterm{def:dv:autocallables:autocallable}{autocallable} has two features that turn hedging into
market flow.

The first is the direction of the flow. The note's value to the investor rises with the underlying, so the issuer, short the note,
holds a long position in the underlying. As the underlying falls towards the \omterm{def:dv:autocallables:protection}{protection barrier}, the note's delta grows, because the
knock-in put the issuer owns gains value quickly there, and the issuer buys. As it rises towards the trigger, the note's delta
shrinks, and the issuer sells. Away from the barrier, then, the issuer's hedging buys dips and sells rallies, and it dampens
volatility. This is the dealer-gamma effect of One Quant Book 1, chapter 26, and the reason issuers are long volatility: they own
the gamma.

The second is the cliff at a daily-monitored barrier. Just above the knock-in level, the note is worth par plus its coupon if the
level holds and much less if it breaks, so its delta is very large. At the knock-in the payoff becomes a plain long position in the
underlying, with a delta of about one. The issuer must sell the difference, at once, into a falling market. Lim and Choi described
this pattern for Korean equity-linked securities: the delta rises continually as the price approaches the knock-in level, and must
be cut as soon as the level is touched.

\begin{example}[The knock-in cliff]\label{ex:dv:autocallables:cliff}
A two-year \omterm{def:dv:autocallables:snowball}{snowball} is observed monthly for knock-out at 103\% from the third month, with a daily knock-in at 75\% and a coupon of
15\% a year. Under \omterm{def:dv:local-volatility:model}{local volatility} it is worth 100.44 (standard error 0.06), and its knock-in probability is 15.6\%. One month
before maturity, at 25\% volatility, with the full 30\% coupon due if the barrier holds and the spot 1\% above the barrier, the issuer's
delta is 7.63 per unit of notional. A 1\% fall knocks the note in, and the delta drops to 1.00. Per 100 million of notes the issuer
sells 497 million of the underlying, measured at the barrier level. With three months left the figure is 254 million, and with a
year left 97 million (\cref{fig:dv:autocallables:cliff}).
\end{example}

\begin{omfigure}
\begin{tikzpicture}
\begin{axis}[width=10.5cm,height=5cm,
  xlabel={underlying (\% of the initial fixing)},ylabel={issuer's delta (per unit of notional)},
  tick label style={font=\small},label style={font=\small},
  xmin=60,xmax=95,ymin=0,ymax=8.5,grid=major,grid style={gray!25},
  legend columns=2,legend style={font=\footnotesize,at={(0.5,-0.3)},anchor=north,draw=none,fill=none,/tikz/every even column/.append style={column sep=8pt}},
  table/col sep=comma]
\addplot[omThm,very thick] table[x=s,y=delta]{figdata/derivatives/18-autocallables/cliff_before.csv};
\addplot[omDef,very thick,dashed] table[x=s,y=delta]{figdata/derivatives/18-autocallables/cliff_after.csv};
\draw[gray,thick] (axis cs:75,0) -- (axis cs:75,8.5);
\node[font=\footnotesize,anchor=west] at (axis cs:76,7.9) {knock-in at 75\%};
\legend{not yet knocked in,after the knock-in}
\end{axis}
\end{tikzpicture}

\omcaption{The issuer's delta on a \omterm{def:dv:autocallables:snowball}{snowball} one month before maturity (knock-in at 75\% monitored daily, 30\% coupon due if the
barrier holds, 25\% volatility). The issuer buys as the underlying falls towards the barrier and, at the knock-in, must sell
most of what it bought. Data: the tutorial.}\label{fig:dv:autocallables:cliff}
\end{omfigure}

Fang's reconstruction of the Chinese \omterm{def:dv:autocallables:snowball}{snowball} episode puts the hedge released across the barrier region at about 0.6 times notional
with twelve months to maturity and about 6 times with one month, the same order as the example. The size of the flow depends on
three things that no single issuer controls: how much notional is outstanding, how concentrated the knock-in levels are, and how
concentrated the maturities are. A dozen issuers selling similar notes in the same months, with barriers at the same round
percentages, create one very large barrier.

\section{Documented market impact}

The structure of the risk is not in doubt. The evidence on its effect on markets is growing, and it should be read with its
limits.

\begin{itemize}
\item \textbf{The 2024 \omterm{def:dv:autocallables:snowball}{snowball} knock-ins.} Yin's case study of the January 2024 ``knock-in storm'' lists concentrated product
maturities and clustered knock-in levels among the causes, and institutional hedging pressure and market volatility among the
effects. Fang finds that basis widening and trading activity in the CSI 500 and CSI 1000 futures, the indices most \omterm{def:dv:autocallables:snowball}{snowballs}
referenced, exceeded those in control index futures around the episode. He describes these as consistency results, not causal
identification.
\item \textbf{Korean equity-linked securities.} Korean studies of the issuer's knock-in delta led to a policy suggestion: size
issuance on the same underlying against its trading volume and the maximum delta the hedge can reach. Losses on notes linked to the
Hang Seng China Enterprises Index are the subject of empirical pricing studies.
\item \textbf{Where the effect is small.} Fang notes that the same monitoring framework, outstanding notional times barrier
concentration times maturity concentration times dealer gamma, transfers to other markets, but the magnitude does not. Japanese
retail notes are more often single-name, with barriers further out of the money.
\end{itemize}

For a desk the practical consequences are three. Know the market's concentration of barriers as well as your own: the flow at a level
is everyone's. Price the cliff: a daily-monitored knock-in near maturity carries \omterm{def:dv:barriers-and-digitals:gap}{gap risk} (chapter 15) and a liquidity cost the model
does not see. And watch the investor's side: the product that pays the highest coupon is the one in which the investor has sold the
most, and the losses when the barrier breaks fall on the investor.

\section{Tutorial: pricing and hedging an autocallable}\label{tut:dv:autocallables:tut}

\begin{tutorial}
\textbf{Goal.} Price a three-year Phoenix under \omterm{def:dv:local-volatility:model}{local volatility}, solve for its fair coupon, compute its \omterm{def:dv:greeks-and-the-hedging-pnl:greeks}{vega} by expiry bucket and its
skew, dividend and correlation sensitivities, and measure the issuer's hedge across a daily knock-in. \textbf{End state:} the four
figures and the numbers of the weekend problem.
\begin{tutsteps}
\item \textbf{The term sheet} as data:
\omcode{code/firm/autocall/firm_autocall.py}{20}{34}{The autocallable term sheet.}\label{lst:dv:autocallables:termsheet}
\item \textbf{Cash flows path by path}: coupons with memory, calls, and redemption at maturity with the knock-in:
\omcode{code/firm/autocall/firm_autocall.py}{63}{95}{Autocallable cash flows on simulated paths.}\label{lst:dv:autocallables:cashflows}
\item \textbf{The hedge across the knock-in}, from the note's value just above the barrier and just after it:
\omcode{code/firm/autocall/firm_autocall.py}{128}{141}{The underlying sold at the knock-in.}\label{lst:dv:autocallables:cliff}
\item \textbf{Run} \texttt{dv\_autocall.sensitivities()}, \texttt{worst\_of()}, \texttt{snowball\_price()}, \texttt{cliff()} and
\texttt{fig\_autocall.py}. The fair coupon uses the fact that the price is linear in the coupon: two simulations with common random
numbers solve for it.
\end{tutsteps}
\textbf{What to change next.} Make the \omterm{def:dv:autocallables:protection}{protection barrier} daily-monitored and recompute the fair coupon; add a step-down trigger (100\%,
95\%, 90\% by year); price the Phoenix under the \omterm{def:dv:stochastic-volatility:heston}{Heston model} of chapter 10 and compare.
\end{tutorial}

\section{Build: the autocallable pricer}\label{bld:dv:autocallables:autocall}

\begin{build}
\bfield{Purpose} The miniature firm's structured-products pricer: term sheets as data, Monte Carlo pricing on any volatility function
and any number of underlyings, fair-coupon solving, and the hedge analytics that the structured-products desk of chapter 19 and the
exotic book of chapter 27 use.
\bfield{Interface} \texttt{TermSheet(obs\_times, trigger, coupon, coupon\_barrier, memory, protection, ki\_daily, snowball\_rate,
first\_call)}; \texttt{simulate(sigma\_fn, s0, t\_end, r, q, n\_paths, seed, steps\_per\_year, corr)}; \texttt{cashflows(ts, times,
perf, r)}; \texttt{price(ts, sigma\_fn, r, q, n\_paths, seed, corr, n\_assets) -> price, se, call\_probs, ki\_prob}; \texttt{snowball\_tail};
\texttt{hedge\_across\_ki(ki, tau, vol, final\_coupon, notional)}.
\bfield{Rules} \omterm{def:dv:greeks-and-the-hedging-pnl:greeks}{Greeks} by bumping with common random numbers; the fair coupon from the price's linearity in the coupon; a daily knock-in
priced with daily paths or the \omterm{def:dv:barriers-and-digitals:shift}{barrier shift}, never as continuous; every price reported with its standard error.
\bfield{Acceptance tests} \texttt{code/firm/autocall/tests/}: without volatility the note is called at once; with a trigger never reached
and no knock-in it pays par plus every coupon; daily monitoring knocks in more often than monitoring at maturity; a worst-of of perfectly
correlated assets is the single-asset note; the \omterm{def:dv:autocallables:snowball}{snowball}'s closed form near maturity matches a simulation; the cliff is largest near
maturity.
\bfield{Stretch} Step-down triggers and worst-of \omterm{def:dv:autocallables:snowball}{snowballs}; local-stochastic volatility (chapter 20); pathwise and likelihood-ratio
\omterm{def:dv:greeks-and-the-hedging-pnl:greeks}{Greeks} for the digital features (chapter 23).
\end{build}

\begin{omsources}
\item P.~Guillaume, ``Autocallable structured products'', \emph{Journal of Derivatives} 22(3) (2015) 73--94.
\item H.~Lim and Y.~Choi, ``Knock-in and stocks market effect due to ELS issuance and hedging'', \emph{Journal of Derivatives and
Quantitative Studies} 23 (2015) 289--321.
\item W.~Fang, ``China's autocallable cycle: dealer hedging flows, market impact, and lessons for Japan's structured-note market'',
SSRN 7346463 (2026).
\item Yin, ``Case study on the knock-in storm of snowball products'', \emph{Journal of Global Economy, Business and Finance} 8(3) (2026)
26--33.
\item Kim, Park and Moon, ``Markov regime-switching in pricing equity-linked securities: an empirical study for losses in HSCEI-linked
products'', \emph{Finance Research Letters} 76 (2025) 106929.
\end{omsources}

\section{Exercises}

\begin{exercise}[$\star$]\label{exo:dv:autocallables:1}
List the options embedded in the chapter's Phoenix and say, for each, whether the investor is long or short.
\end{exercise}

\begin{exercise}[$\star$]\label{exo:dv:autocallables:2}
Why is the fair coupon under flat volatility lower than under \omterm{def:dv:local-volatility:model}{local volatility}?
\end{exercise}

\begin{exercise}[$\star$]\label{exo:dv:autocallables:3}
The price is linear in the coupon. Using the chapter's numbers (95.99 without coupons, 2.60 per point), what quarterly coupon prices the
note at 98, leaving the issuer two points of margin?
\end{exercise}

\begin{exercise}[$\star\star$]\label{exo:dv:autocallables:4}
Explain why an issuer of \omterm{def:dv:autocallables:autocallable}{autocallables} is long dividends, and what it does about it.
\end{exercise}

\begin{exercise}[$\star\star$]\label{exo:dv:autocallables:5}
Explain the direction of the issuer's hedging flow as the underlying falls from 100\% towards a daily knock-in, and at the knock-in.
\end{exercise}

\begin{exercise}[$\star\star$]\label{exo:dv:autocallables:6}
Why does the cliff grow as maturity approaches?
\end{exercise}

\begin{exercise}[$\star\star\star$]\label{exo:dv:autocallables:7}
\emph{Coding.} Make the Phoenix's \omterm{def:dv:autocallables:protection}{protection barrier} daily-monitored and recompute the fair coupon under \omterm{def:dv:local-volatility:model}{local volatility}. Explain the
change.
\end{exercise}

\begin{exercise}[$\star\star\star$]\label{exo:dv:autocallables:8}
\emph{Find the flaw.} ``Our \omterm{def:dv:autocallables:autocallable}{autocallable} book is delta-hedged and vega-hedged, so a fall through the knock-in level of our notes is a
non-event for the desk.''
\end{exercise}

\section{Problem: The Knock-in Cliff}

\begin{problem}[{Weekend problem --- how much the issuer sells at the barrier}]\label{pb:dv:autocallables:1}
An issuer has sold \omterm{def:dv:autocallables:snowball}{snowballs} on an index: monthly knock-out at 103\% from the third month, daily knock-in at 75\%, 15\% a year, two
years, 100 million of notional.

\medskip\noindent\textbf{Part I --- The product.}
\begin{enumerate}
\item Price the \omterm{def:dv:autocallables:snowball}{snowball} under chapter 9's \omterm{def:dv:local-volatility:model}{local volatility} and give its knock-in probability.
\item Which options has the investor sold?
\item What does the investor receive in each of the three outcomes (knock-out, neither, knock-in)?
\item Why is the knock-in monitored daily a harder risk than one observed at maturity?
\item Compare with the chapter's Phoenix: which pays more, and why?
\end{enumerate}

\medskip\noindent\textbf{Part II --- The issuer's delta.}
\begin{enumerate}[resume]
\item One month before maturity, at 25\% volatility, spot 1\% above the barrier, give the issuer's delta per unit of notional.
\item Give the delta after a 1\% fall knocks the note in.
\item How much of the underlying does the issuer sell, per 100 million of notes?
\item Give the same with three months and with a year left.
\item What did the issuer do as the spot fell from 100\% to 76\%?
\end{enumerate}

\medskip\noindent\textbf{Part III --- The market.}
\begin{enumerate}[resume]
\item If ten issuers hold similar notes with the same barrier, what does the market see?
\item What order of magnitude does Fang report for the hedge released across the barrier region?
\item Why are futures on the index where the flow shows first?
\item What policy did Lim and Choi suggest?
\item Which data would a regulator need to monitor this risk?
\end{enumerate}

\medskip\noindent\textbf{Part IV --- Judgement.}
\begin{enumerate}[resume]
\item How would you price the cliff's liquidity cost into the note?
\item How would you hedge ahead of the barrier to reduce the sale at the knock-in?
\item Who bears the loss when the barrier breaks?
\item State the \emph{named result}: the quantity of underlying the issuer sells per 1\% fall when spot is 1\% above the \omterm{def:dv:autocallables:protection}{protection
barrier} one month before maturity, per 100 million of notes.
\item In one sentence: what makes a knock-in barrier a market event rather than a single trade's event?
\end{enumerate}
\end{problem}

\section{Interview questions}

\begin{interviewq}[$\star$ \iqroles{trader, bank}]\label{iq:dv:autocallables:1}
Decompose an \omterm{def:dv:autocallables:autocallable}{autocallable} into options. Why can it pay a high coupon?
\end{interviewq}

\begin{interviewq}[$\star\star$ \iqroles{researcher}]\label{iq:dv:autocallables:2}
Which model would you price an \omterm{def:dv:autocallables:autocallable}{autocallable} with, and what does a flat volatility get wrong?
\end{interviewq}

\begin{interviewq}[$\star\star$ \iqroles{trader}]\label{iq:dv:autocallables:3}
What is an \omterm{def:dv:autocallables:autocallable}{autocallable} issuer's \omterm{def:dv:greeks-and-the-hedging-pnl:greeks}{vega}, skew and dividend exposure, and how is it recycled?
\end{interviewq}

\begin{interviewq}[$\star\star$ \iqroles{developer}]\label{iq:dv:autocallables:4}
Your Monte Carlo \omterm{def:dv:greeks-and-the-hedging-pnl:greeks}{Greeks} for an \omterm{def:dv:autocallables:autocallable}{autocallable} are noisy. What do you do?
\end{interviewq}

\begin{interviewq}[$\star\star$ \iqroles{risk}]\label{iq:dv:autocallables:5}
How would you measure the market-wide hedging flow that a cluster of knock-in barriers can create?
\end{interviewq}

\begin{interviewq}[$\star\star\star$ \iqroles{trader, risk}]\label{iq:dv:autocallables:6}
The index is 2\% above a level where a large share of the market's notes knock in, one month before many of them mature. What do you
do?
\end{interviewq}
